%! Summary Statistics for a Quantitative Variable — Unit 1, Lesson 1.5 — Experience & Formalize
%! (Activity · QuickNotes · Application) (Answer Key)
% EFFL component, Math Medic sizing: 12pt body, OPEN white space after every question.
% Page budget: 3pp total — Activity <=2pp · QuickNotes 1/2pp · Application 1/2-1pp.
% ANSWERSPACE BUDGET: 11.4cm across 7 sub-questions (1.4 + 2.0 + 1.6 + 1.4 + 1.8 + 1.8 + 1.6),
% in line with the model lesson (unit01/lesson03, 10.8cm across 8). Every key answer is written
% to FIT its reserved height at ~70 characters per line — do not raise a height without
% shortening the answer, and never raise one in only one of the two files.
% Check Your Understanding is NOT in this component: those items were moved to the end of
% the packet and are now the lesson's SCORED homework (see ../homework/).
%
% SPOILER RULE: the Activity never says statistic, quartile, Q1, Q3, interquartile range, IQR,
% percentile, range (as one number), standard deviation, variance, resistant, or nonresistant.
% It says "the average", "the value in the middle", "one number for how spread out". Shape,
% skew, centre, spread, and outlier WERE formalized in Lesson 1.4 and are prior knowledge.
% Everything from the QuickNotes box onward is post-debrief, so the Application uses the full
% AP vocabulary deliberately.
%
% NO SKETCH QUESTIONS (hard constraint): every dotplot is pre-drawn. Students "construct" only
% by filling the five-number-summary table in the Application, never by drawing an axis.
%
% THE DATA
%   Problem 1 — nine weekly paychecks:  650 700 700 750 750 800 800 850 3000
%     sum 9000 -> mean $1,000 ; median (5th of 9) $750.
%     Replace the 3000 by 750: sum 6750 -> mean $750 (a $250 fall); median STILL $750.
%     That is the crux (EK UNC-1.K.2: mean nonresistant, median resistant).
%   Problem 2 — twelve delivery times per route, all three routes:
%     A: 27 28 28 29 29 30 30 31 31 32 32 33   sum 360, mean 30, median 30, max-min = 6
%     B: 18 20 22 24 26 29 31 34 36 38 40 42   sum 360, mean 30, median 30, max-min = 24
%     C: 27 28 28 29 29 30 30 31 31 32 32 51   median 30, max-min = 24  <- same as B, from ONE
%     value. That collapse is the seed for the IQR (EK UNC-1.J.2).
% GENERATED FROM experience/main.tex — the blank with answers filled in, so every \answerspace
% height and \boxguard is identical to it. NO teachernote here — that lives in the lesson plan.
% \ans{} IS TEXT MODE: never inside $...$; a filled slot inside a formula is written
% {\color{keyred}\mathbf{...}} instead — see the standard-deviation line in the QuickNotes.
\documentclass[12pt]{article}
\usepackage{apstats-article}
\usepackage{apstats-key}   % pulls in -boxes; do NOT also load -boxes
\newcommand{\answerspace}[2]{\par\nopagebreak\noindent\begin{minipage}[t][#1][t]{\linewidth}%
  \color{keyred}\bfseries #2\end{minipage}\par}

\begin{document}

\pageheader{Unit 1, Lesson 1.5}{Experience \& Formalize: Nine Paychecks --- Answer Key}
% NAMESTRIP: no \namedateperiod here — mirrors the blank.

\vspace{0.02in}

% ── Part 1: the Activity ─────────────────────────────────────────────────────
\begin{headlinebox}{sky}
\textbf{Nine Paychecks.} A small landscaping company has nine people on the payroll --- eight
crew members and the owner. Below is last week's take-home pay, and after it, delivery times
from a courier company. Work both problems with your group using only what you already know.
\end{headlinebox}

\vspace{0.04in}

\begin{scenariobox}[1. Nine Paychecks]{navy}
Last week's take-home pay, in dollars, for all nine people at the company:

\vspace{3pt}
\begin{center}
800 \quad 3000 \quad 650 \quad 750 \quad 700 \quad 850 \quad 700 \quad 800 \quad 750
\end{center}
\vspace{2pt}

\begin{center}
\begin{tikzpicture}[scale=0.75]
  \draw[->, thick] (-0.3,0) -- (13.8,0);
  \foreach \x/\lab in {0/600, 2.16/1000, 4.32/1400, 6.48/1800, 8.64/2200, 10.8/2600, 12.96/3000} {
    \draw (\x,-0.08) -- (\x,0.08); \node[below, font=\scriptsize] at (\x,-0.11) {\lab};}
  \foreach \x/\n in {0.27/1, 0.54/2, 0.81/2, 1.08/2, 1.35/1, 12.96/1} {
    \foreach \k in {1,...,\n} { \fill[navy] (\x,{0.2*\k}) circle (0.07); }}
  \node[font=\scriptsize] at (6.5,-0.60) {Weekly take-home pay (dollars)};
\end{tikzpicture}
\end{center}
\vspace{1pt}

\normalsize
\begin{enumerate}[label=\textbf{\alph*.}, leftmargin=1.8em, itemsep=4pt]
  \item Find the \textbf{average} of the nine paychecks. \ans{\$1{,}000}

        Now put the nine in order and give the value in the \textbf{middle}. \ans{\$750 (the 5th)}
  \item A crew member says, \emph{``The typical weekly pay here is a thousand dollars.''} Is she
        quoting a number that is actually in this list? How many of the nine earn that much?
        \answerspace{1.4cm}{No --- nobody here is paid \$1{,}000. Only the owner takes home that
        much; the other eight all earn under \$900.}
  \item \textbf{Cross out the owner's \$3{,}000} and replace it with \textbf{\$750}. Recompute
        both answers from part (a). Say which one barely moved, by how much each moved, and
        \emph{why the arithmetic makes that happen}.
        \answerspace{2.0cm}{New average \$750 (\$6{,}750 over 9) --- a fall of \$250. The middle
        value is still \$750. The average adds all nine numbers, so one huge paycheck lifts the
        total; the middle value only asks which paycheck is fifth.}
  \item Which of your two numbers would you put in a \textbf{job advertisement}, and which in a
        \textbf{union negotiation}? Defend both choices.
        \answerspace{1.6cm}{The advertisement: \$1{,}000, the bigger number and honestly
        computed. The negotiation: \$750, because eight of the nine earn near it and nobody
        earns \$1{,}000. Both are correct --- different questions.}
\end{enumerate}
\end{scenariobox}

\boxguard[16]
\begin{scenariobox}[2. Two Delivery Routes]{navy}
A courier company times how long each of twelve packages takes to reach its door, on two
different routes. Each dot is one package.

\vspace{3pt}
\begin{center}
\begin{tikzpicture}[scale=0.7]
  % ── left: Route A ──
  \draw[->, thick] (-0.2,0) -- (6.9,0);
  \foreach \x/\lab in {0.8/20, 2.4/30, 4.0/40, 5.6/50} {
    \draw (\x,-0.08) -- (\x,0.08); \node[below, font=\scriptsize] at (\x,-0.11) {\lab};}
  \foreach \x/\n in {1.92/1, 2.08/2, 2.24/2, 2.40/2, 2.56/2, 2.72/2, 2.88/1} {
    \foreach \k in {1,...,\n} { \fill[navy] (\x,{0.18*\k}) circle (0.065); }}
  \node[font=\scriptsize] at (3.2,-0.56) {Delivery time (minutes)};
  \node[font=\bfseries\small] at (3.2,1.0) {Route A};
  % ── right: Route B ──
  \begin{scope}[xshift=7.6cm]
    \draw[->, thick] (-0.2,0) -- (6.9,0);
    \foreach \x/\lab in {0.8/20, 2.4/30, 4.0/40, 5.6/50} {
      \draw (\x,-0.08) -- (\x,0.08); \node[below, font=\scriptsize] at (\x,-0.11) {\lab};}
    \foreach \x/\n in {0.48/1, 0.80/1, 1.12/1, 1.44/1, 1.76/1, 2.24/1, 2.56/1, 3.04/1,
                       3.36/1, 3.68/1, 4.00/1, 4.32/1} {
      \foreach \k in {1,...,\n} { \fill[goldacc] (\x,{0.18*\k}) circle (0.065); }}
    \node[font=\scriptsize] at (3.2,-0.56) {Delivery time (minutes)};
    \node[font=\bfseries\small] at (3.2,1.0) {Route B};
  \end{scope}
\end{tikzpicture}
\end{center}
\vspace{1pt}

\normalsize
\begin{enumerate}[label=\textbf{\alph*.}, leftmargin=1.8em, itemsep=4pt]
  \item \textbf{Each route's twelve times add to 360 minutes.} With twelve values none sits
        alone in the middle, so take the value halfway between the sixth and the seventh.

        \vspace{3pt}
        Route A --- average: \ans{30 min} \quad middle: \ans{30 min}

        Route B --- average: \ans{30 min} \quad middle: \ans{30 min}

        \vspace{3pt}
        Do the two routes have the same centre? \ans{Yes --- both}
  \item Those summaries cannot tell these routes apart, and anybody can see they are not the
        same. \textbf{Invent a number that would.} Say how to compute it and give both values.
        \answerspace{1.4cm}{Largest minus smallest. Route A: $33-27=6$ min. Route B:
        $42-18=24$ min. One number, and it separates them at once.}
  \item Here is a \textbf{third} route from the same company.

        \vspace{2pt}
        \begin{center}
        \begin{tikzpicture}[scale=0.7]
          \draw[->, thick] (-0.2,0) -- (6.9,0);
          \foreach \x/\lab in {0.8/20, 2.4/30, 4.0/40, 5.6/50} {
            \draw (\x,-0.08) -- (\x,0.08); \node[below, font=\scriptsize] at (\x,-0.11) {\lab};}
          \foreach \x/\n in {1.92/1, 2.08/2, 2.24/2, 2.40/2, 2.56/2, 2.72/2, 5.76/1} {
            \foreach \k in {1,...,\n} { \fill[navy] (\x,{0.18*\k}) circle (0.065); }}
          \node[font=\scriptsize] at (3.2,-0.56) {Delivery time (minutes)};
          \node[font=\bfseries\small] at (3.2,1.0) {Route C};
        \end{tikzpicture}
        \end{center}
        \vspace{1pt}

        Work out your number from part (b) for Route C and compare it with Route B's. Then say
        what your number is claiming about Route C that simply is not true.
        \answerspace{1.8cm}{Route C: $51-27=24$ min --- Route B's number exactly. It claims C is
        as unpredictable as B, but eleven of C's twelve packages arrived between 27 and 32.
        Largest-minus-smallest sees only the two ends.}
  \item \textbf{Fix it.} Describe a way of measuring how spread out a set of times is that the
        one late package \emph{cannot} hijack --- clearly enough that someone else could carry
        it out.
        \answerspace{1.8cm}{Ignore the far ends; measure the middle. Order the times, take the
        value a quarter of the way up and the value three quarters up, and subtract. One late
        package cannot touch that middle stretch.}
\end{enumerate}
\end{scenariobox}

% ── Part 2: QuickNotes (the debrief fills this) — target: half a page ────────
\boxguard[14]
\begin{tcolorbox}[enhanced, colback=sky, colframe=navy, boxrule=0.8pt, arc=2mm,
    left=3mm, right=3mm, top=1.5mm, bottom=1.5mm, breakable,
    fonttitle=\bfseries\small\color{white}, title={QuickNotes: Summarizing a Quantitative Variable},
    attach boxed title to top left={yshift=-2mm, xshift=4mm},
    boxed title style={colback=navy, colframe=navy, arc=1.5mm}]
\small
% Display skips are tightened so the standard-deviation formula does not cost half a page.
\setlength{\abovedisplayskip}{3pt}\setlength{\belowdisplayskip}{3pt}
\setlength{\abovedisplayshortskip}{3pt}\setlength{\belowdisplayshortskip}{3pt}
\begin{itemize}[leftmargin=1.1em, itemsep=3pt]
  \item A \textbf{statistic} is a numerical summary of \ans{sample} data.
        \textbf{Mean} ($\bar x$): add every value and divide by \ans{the number of values}.
        \textbf{Median}: the \ans{middle} value of the \emph{ordered} data --- with an even
        count, the \ans{average} of the two middle values.
  \item The \textbf{$p$th percentile} is the value with $p\%$ of the data
        \ans{less than or equal to} it.
        \textbf{$Q_1$} and \textbf{$Q_3$} are the medians of the lower and upper halves;
        between them sits the middle \ans{50}\% of the data.
  \item \textbf{Range} $=$ \ans{maximum} $-$ \ans{minimum}. \quad
        \textbf{IQR} $=$ \ans{$Q_3$} $-$ \ans{$Q_1$}. \quad Each is \emph{one number}, not
        an interval.
  \item \textbf{Standard deviation} ($s_x$): roughly the \ans{typical} distance of a value from
        the mean. We get it from \ans{technology}; its square $s_x^2$ is the \ans{variance}.
        \[ s_x \;=\; \sqrt{\frac{1}{n-{\color{keyred}\mathbf{1}}}\sum_{i=1}^{n}\left(x_i-\bar x\right)^2} \]
  \item \textbf{Resistant} (one extreme value barely moves it): \ans{median} and
        \ans{IQR}. \quad \textbf{Nonresistant} (one extreme value drags it):
        \ans{mean}, \ans{standard deviation}, and \ans{range}.
  \item \textbf{Two outlier rules in this course:} more than $1.5\times\text{IQR}$ above $Q_3$
        or below $Q_1$, or 2 or more standard deviations from the mean. The
        $1.5\times\text{IQR}$ rule goes to work \ans{next} lesson.
\end{itemize}
\end{tcolorbox}

% ── Part 3: the Application (worked together, ~6 min) ────────────────────────
\boxguard[14]
\begin{notesbox}{Application: The Quiz Scores}
We will do this one together. Eleven students took a 20-point quiz. Their scores, in order:

\vspace{2pt}
\begin{center}
12 \quad 13 \quad 14 \quad 15 \quad 15 \quad 16 \quad 17 \quad 18 \quad 18 \quad 19 \quad 19
\end{center}
\vspace{1pt}

\begin{center}
\begin{tikzpicture}[scale=0.8]
  \draw[->, thick] (-0.2,0) -- (7.5,0);
  \foreach \x/\lab in {0/10, 1.4/12, 2.8/14, 4.2/16, 5.6/18, 7.0/20} {
    \draw (\x,-0.08) -- (\x,0.08); \node[below, font=\scriptsize] at (\x,-0.11) {\lab};}
  \foreach \x/\n in {1.4/1, 2.1/1, 2.8/1, 3.5/2, 4.2/1, 4.9/1, 5.6/2, 6.3/2} {
    \foreach \k in {1,...,\n} { \fill[navy] (\x,{0.22*\k}) circle (0.075); }}
  \node[font=\scriptsize] at (3.6,-0.60) {Quiz score (points out of 20)};
\end{tikzpicture}
\end{center}
\vspace{1pt}

\begin{enumerate}[label=\textbf{\alph*.}, leftmargin=1.8em, itemsep=4pt]
  \item Find the median, then split the ordered list at it --- five scores below, five above.

        \vspace{2pt}
        \renewcommand{\arraystretch}{1.4}
        \begin{tabularx}{\linewidth}{@{} Y Y Y Y Y @{}}
        \rowcolor{sky}
        \textbf{Minimum} & $\mathbf{Q_1}$ & \textbf{Median} & $\mathbf{Q_3}$ & \textbf{Maximum} \\
        \hline
        \ans{12} & \ans{14} & \ans{16} & \ans{18} & \ans{19} \\
        \end{tabularx}
  \item Compute the \textbf{range} and the \textbf{IQR} for these eleven scores.

        \begin{work}
        \text{range} &= \text{maximum} - \text{minimum} = 19 - 12 = 7 \text{ points} \\
        \text{IQR} &= Q_3 - Q_1 = 18 - 14 = 4 \text{ points}
        \end{work}
  \item \textbf{One value changes.} The lowest score was entered wrong: the \textbf{12 was
        really a 1}. Of \emph{mean, median, range, IQR}, which change and which do not? State
        the rule that decides it.
        \answerspace{1.6cm}{Mean $16\to15$ and range $7\to18$ change; median stays 16 and IQR
        stays 4. Mean and range are built from the extreme values (nonresistant); median and
        IQR use only middle positions (resistant).}
\end{enumerate}
\end{notesbox}

\end{document}
